\subsection{Polstellen: Was nach dem Kürzen übrig bleibt}
\label{subsec:real_functions_poles}
Nicht jede ausgeschlossene Eingabe lässt sich ergänzen.
Die Vielfachheiten der gemeinsamen Faktoren entscheiden darüber,
ob eine Lücke hebbar ist oder die Werte unbegrenzt werden.
\par\noindent\begin{minipage}{\linewidth}
\begin{DefinitionBox}{Polstelle}{real_functions_pole}
  Eine ausgeschlossene Stelle $a$ einer gebrochen-rationalen Funktion
  heißt \emph{Polstelle} \EnglishTerm{pole}, wenn
  $\lim_{x\to a}|f(x)|=+\infty$ gilt.
  Die Gerade $x=a$ heißt dann \emph{senkrechte Asymptote}
  \EnglishTerm{vertical asymptote}.
\end{DefinitionBox}
\begin{ExampleBox}[before skip=0.3\baselineskip]{Gemeinsame Nullstellen reichen nicht aus}{real_functions_common_zeros}
  Bei allen drei Termen werden Zähler und Nenner für $x\to1$ null.
  Kürzen Sie und entscheiden Sie, welches Verhalten bleibt.
  \[
    f(x)=\frac{x-1}{x-1},\qquad
    g(x)=\frac{(x-1)^2}{x-1},\qquad
    h(x)=\frac{x-1}{(x-1)^2}.
  \]
  \begin{ClozeArea}[34]{10}
    \[
      \begin{array}{l@{\quad}l}
        \dfrac{x-1}{x-1}=1
          &\text{hebbare Lücke, Grenzwert }1,\\[1.1em]
        \dfrac{(x-1)^2}{x-1}=x-1
          &\text{hebbare Lücke, Grenzwert }0,\\[1.1em]
        \dfrac{x-1}{(x-1)^2}=\dfrac1{x-1}
          &\text{Polstelle}.
      \end{array}
      \qquad(x\ne1)
    \]
    Im dritten Fall werden die Werte von $h$ bei der Annäherung an $1$
    von links beliebig stark negativ, von rechts beliebig groß positiv.
  \end{ClozeArea}
\end{ExampleBox}
\end{minipage}


\begin{KeyPointBox}{Vielfachheiten vergleichen}{real_functions_pole_order}
  An einer Nennernullstelle $a$ schreiben wir mit $r(a),s(a)\ne0$
  \[
    f(x)=\frac{(x-a)^m r(x)}{(x-a)^n s(x)}
        =(x-a)^{m-n}\frac{r(x)}{s(x)},\qquad n\ge1,\ m\ge0.
  \]
  Für $m\ge n$ ist die Lücke hebbar. Für $m<n$ bleibt eine Polstelle
  der Ordnung $n-m$:
  \begin{itemize}
    \item ungerade Ordnung: entgegengesetzte Vorzeichen auf beiden Seiten;
    \item gerade Ordnung: dasselbe Vorzeichen auf beiden Seiten.
  \end{itemize}
  Das konkrete Vorzeichen ergibt sich zusätzlich aus $r(a)/s(a)$.
\end{KeyPointBox}
% Handschrift: PDF-Seite 24; gemeinsame Nullstellen nach Vielfachheit präzisiert.
\begin{center}
  \CachedTikzFigure{content/resources/chapters/real_functions/tikz/functions_rational_decision.tex}
\end{center}
\begin{ExampleBox}{Eine Nullstelle ohne Definitionslücke und eine Polstelle}{real_functions_rational_zero_and_pole}
  Vergleichen Sie die Zähler- und Nennernullstellen der beiden Funktionen.
  Bestimmen Sie anschließend Nullstellen, Definitionsmengen und Polstellen.
  \begin{enumerate}[label=\alph*)]
    \item $f(x)=\dfrac{x}{x^2+1}$: Der Nenner ist stets positiv.
      \[
        D_f=\Cloze{2cm}{\R},\qquad
        f(x)\overset{!}{=}0\ \Longleftrightarrow\ x=\Cloze{1cm}{0}.
      \]
    \item $g(x)=\dfrac1{x-1}$: Der Zähler wird nie null.
      \[
        D_g=\Cloze{3cm}{\R\setminus\{1\}}.
      \]
      Bei der Annäherung an $1$ werden die Werte von links
      \Cloze{5cm}{beliebig stark negativ}, von rechts
      \Cloze{5cm}{beliebig groß positiv}.
      Es gibt keine Nullstelle; $x=1$ ist eine Polstelle.
  \end{enumerate}
  \begin{center}
    \CachedTikzFigure{content/resources/chapters/real_functions/tikz/functions_rational_zero_and_pole.tex}
  \end{center}
  Die hebbare Lücke von $(x^2-1)/(x+1)$ wurde zuvor in
  Beispiel~\ref{ex:real_functions_cancelled_hole} untersucht.
\end{ExampleBox}

\clearpage
\subsection{Asymptoten im Unendlichen}
\label{subsec:real_functions_asymptotes}
Ein Graph kann sich für große Eingaben einer Geraden annähern.
Dabei muss der Unterschied der Funktionswerte gegen null gehen --
ein bloß ähnlicher Verlauf genügt nicht.
\par\noindent\begin{minipage}{\linewidth}
\begin{DefinitionBox}{Waagerechte und schiefe Asymptote}{real_functions_asymptote}
  Die Gerade $y=mx+b$ ist eine \emph{Asymptote} \EnglishTerm{asymptote}
  für $x\to+\infty$ beziehungsweise $x\to-\infty$, wenn
  \[
    \lim_{x\to\pm\infty}\bigl(f(x)-(mx+b)\bigr)=0
  \]
  für die betreffende Richtung gilt.
  Für $m=0$ heißt sie \emph{waagerecht}
  \EnglishTerm{horizontal}, für $m\ne0$ \emph{schief}
  \EnglishTerm{oblique}.
\end{DefinitionBox}
\begin{ExampleBox}[before skip=0.3\baselineskip]{Die Polynomdivision zeigt die Asymptote}{real_functions_slant_asymptote}
  Wir untersuchen $f(x)=(x^2+1)/(x-1)$ für große positive und negative $x$.
  \begin{ClozeArea}[34]{6}
    Die Polynomdivision zerlegt den Term in eine Gerade und einen Rest:
    \[
      f(x)=x+1+\frac2{x-1}.
    \]
    Weil $\lim_{x\to\pm\infty}2/(x-1)=0$, ist $y=x+1$
    auf beiden Seiten eine schiefe Asymptote.
    Zusätzlich ist $x=1$ eine senkrechte Asymptote.
  \end{ClozeArea}
  \begin{center}
    \CachedTikzFigure{content/resources/chapters/real_functions/tikz/functions_slant_asymptote.tex}
  \end{center}
\end{ExampleBox}
\end{minipage}


\begin{TheoremBox}{Gradvergleich für gebrochen-rationale Funktionen}{real_functions_degree_comparison}
  Seien $m=\deg p$, $n=\deg q$ und $a_m,b_n$ die Leitkoeffizienten.
  Für $f=p/q$ mit $p\ne0$ gilt:
  \begin{center}
    \renewcommand{\arraystretch}{1.5}
    \begin{tabular}{@{}ll@{}}
      \toprule
      Gradvergleich & Verhalten für $x\to\pm\infty$\\ \midrule
      $m<n$ & $\lim_{x\to\pm\infty}f(x)=0$\\
      $m=n$ & $\lim_{x\to\pm\infty}f(x)=a_m/b_n$\\
      $m=n+1$ & schiefe Asymptote aus Polynomdivision\\
      $m>n+1$ & Annäherung an ein Polynom vom Grad $m-n$\\
      \bottomrule
    \end{tabular}
  \end{center}
  Allgemein liefert die Division $p=hq+r$ mit $\deg r<n$ die Zerlegung
  $f=h+r/q$. Wegen $\lim_{x\to\pm\infty}r(x)/q(x)=0$ beschreibt $h$
  das Verhalten im Unendlichen.
\end{TheoremBox}
% Handschrift: PDF-Seiten 27–29 (Herleitung und sämtliche Gradvergleichsfälle).
Durch Ausklammern der höchsten Potenzen wird der Gradvergleich sichtbar:
\[
  \begin{gathered}
    f(x)=x^{m-n}
    \frac{a_m+a_{m-1}/x+\dots+a_0/x^m}
         {b_n+b_{n-1}/x+\dots+b_0/x^n},\\[0.6em]
    \lim_{x\to\pm\infty}
      \frac{a_m+a_{m-1}/x+\dots+a_0/x^m}
           {b_n+b_{n-1}/x+\dots+b_0/x^n}=\frac{a_m}{b_n}.
  \end{gathered}
\]
Für $m>n$ wächst $|f(x)|$ daher gegen $\infty$; das Vorzeichen ergibt
sich aus $a_m/b_n$ und der Parität von $m-n$.

\begin{ExampleBox}{Zwei waagerechte Asymptoten}{real_functions_horizontal_asymptote_examples}
  Dividieren Sie Zähler und Nenner durch die höchste Nennerpotenz.
  Bestimmen Sie so die Asymptote für beide Richtungen im Unendlichen.
  \begin{ClozeArea}[34]{7}
    \[
      \begin{aligned}
        \lim_{x\to\pm\infty}\frac{x}{x^2+1}
          &=\lim_{x\to\pm\infty}\frac{1/x}{1+1/x^2}=0,\\[0.8em]
        \lim_{x\to\pm\infty}\frac{4x^2+2x+1}{x^2+1}
          &=\lim_{x\to\pm\infty}\frac{4+2/x+1/x^2}{1+1/x^2}=4.
      \end{aligned}
    \]
    Die waagerechten Asymptoten sind $y=0$ beziehungsweise $y=4$.
  \end{ClozeArea}
  \begin{center}
    \CachedTikzFigure{content/resources/chapters/real_functions/tikz/functions_horizontal_asymptotes.tex}
  \end{center}
\end{ExampleBox}
\begin{ExampleBox}{Annäherung an eine Parabel}{real_functions_polynomial_asymptote}
  Für $f(x)=(x^4-2x^2+1)/(x^2+1)$ ist der Zählergrad um $2$ größer.
  Bestimmen Sie durch Polynomdivision das Polynom, dem sich der Graph nähert.
  \begin{ClozeArea}[34]{9}
    \[
      \begin{array}{r@{\;}l}
        (x^4-2x^2+1):(x^2+1)&=x^2-3+\dfrac4{x^2+1}\\[0.8em]
        \underline{-(x^4+x^2)}&\\[0.5em]
        -3x^2+1&\\[0.5em]
        \underline{-(-3x^2-3)}&\\[0.5em]
        4&
      \end{array}
    \]
    Mit $A(x)=x^2-3$ gilt
    $\lim_{x\to\pm\infty}\bigl(f(x)-A(x)\bigr)=
      \lim_{x\to\pm\infty}4/(x^2+1)=0$.
    Außerdem ist $\lim_{x\to\pm\infty}f(x)=+\infty$.
  \end{ClozeArea}
  \begin{center}
    \CachedTikzFigure{content/resources/chapters/real_functions/tikz/functions_polynomial_asymptote.tex}
  \end{center}
\end{ExampleBox}
\begin{RemarkBox}
  Für $m>n+1$ ist der Graph von $h$ eine polynomial beschriebene
  Annäherungskurve, keine schiefe Gerade. Eine Asymptote darf außerdem
  vom Funktionsgraphen geschnitten werden.
\end{RemarkBox}
